\section{Appendix: Full Proof}
\label{sec-appendix}

In this section, we establish key properties of the selective view materialization algorithm, including the minimal form uniqueness, the termination guarantee, and the completeness, which together form the correctness of \proto algorithm.

\medskip \noindent {\bf Minimal Form Uniqueness.}
We first define the minimal form of a materialization plan, and then prove the uniqueness of the minimal form for plans generated by our algorithm.
%We denote a valid view materialization plan as a valid set for conciseness.

\begin{deff}
A relation set $S'$ is a minimal form of another relation set $S$ if and only if:
\begin{enumerate}
    \item 
    $S' \subseteq S$, 
    \item
    $S'$ is a minimal view materialization plan (Definition~\ref{def:min-plan}).
\end{enumerate}
\end{deff}

\begin{lmm}\label{lemma:uniqness}
    A valid materialization plan $S$ generated following the base minimal plan has a unique minimal form $S'$.
\end{lmm}

\begin{proof}
    We proceed by establishing contradictions. 
    Assume that a valid materialization plan S has two distinct minimal forms, denoted as X and Y. 
    To demonstrate the impossibility of such a scenario, we break down the proof by examining various overlapping cases of X and Y.

    First, suppose $X \subset Y$, and let $D=Y \setminus X$. 
    Since $X$ is a valid relation set, we can obtain another valid set by removing relations in $D$ from $Y$, which contradicts the assumption that $Y$ is in a minimal form.
    Similarly, the case where $Y \subset X$ also leads to a contradiction.

    The remaining case is where sets X and Y have at least one element that is unique to each set.
    Let $A$ be the set of unique elements in $X$ ($X \setminus Y$), and let $B$ be the set of unique elements in $Y$ ($ Y \setminus X$). We have $A \neq \emptyset$ and $B \neq \emptyset$.

    Next, we want to show that for each relation $a$ in $A$, 
    the execution of queries to $a$ depends on at least one relation in $B$.
    To see this, consider the fact that $S$ is a valid materialization plan,
    and that $Y$ is a minimal form of $S$.
    Consequently, relational queries to $a$ should be executable given $Y$.
    Moreover, if there exists a relation $a$ in $A$ such that $a$ is executable solely 
    using relations in $X \cap Y$, then we can obtain another valid set by removing $a$ from $X$,
    which contradicts the assumption that $X$ is in a minimal form.
    Therefore, queries to each relation $a$ in $A$ must rely on at least one relation in $B$.
    
    Similarly, queries to each relation $b$ in $B$ must also rely on at least one relation in $A$.
    This yields mutual dependencies between relations in set $A$ and $B$, contradicting the assumption that any contract has no circular dependencies
    among relations.

    Therefore, we successfully establish the uniqueness of the minimal form of a given valid
    view materialization plan.
\end{proof}

\medskip\noindent{\bf Termination.} We next briefly discuss the termination property.

\begin{lmm}\label{lemma:uniqness}
    The enumeration algorithm is guaranteed to terminate.
\end{lmm}

\begin{proof}
    The enumeration algorithm (Algorithm~\ref{alg:enumerate}) traverses the replacement graph. Given that both the contract and the replacement graph should have no cycle, the traversal is guaranteed to terminate at nodes without any outgoing edges.
\end{proof}
    
    % In theory, there is no circle in the replacement graph after we remove all the transaction nodes together with their incoming edges. Otherwise, it means that there is at least one cycle composed of non-transaction nodes. And since for non-transaction rules (those not containing "recv\_"), they will be updated immediately, instead of triggered by a function call, at the point when a body relation changes, which will lead to unterminated updation when a relation in the circle of non-transaction nodes changes, indicating an invalid datalog. 
% 
    % Therefore, considering that there is no circle in the replacement graph, we can further use dynamic programming to prove that there are limited number of alternative min plans. Let $N(v)$ denote the upper bound of the number of alternative choices for node $v$, considering both its direct body relations and secondary or further body relations. $\forall v$, we have $N(v) = 1 + \prod N(u)$ where $u$ belongs to the body nodes of $v$ and 1 refers to itself when we do not replace $v$. The base case will be the nodes without incoming edges, and for those $v$ we have $N(v) =1$. Besides, we have $maximum \ number \ of \ min \ sets \leq \prod N(v)$ where $v$ belongs to the base min set. Since there are limited nodes and edges in the replacement graph without circles, the replacement operations can be conducted a finite number of times to reach the end and the $maximum \ number \ of \ min \ sets$ is a finite number, which indicates that this algorithm terminates.
% \end{proof}

\medskip \noindent{\bf Completeness.}
We first define the $Replace$ and $GetMinimal$ functions and then establish the completeness of our proposed algorithm.

We denote the $Replace$ function on a set of view materialization plans $\Sigma$ as follows:
\[ Replace(\Sigma) \triangleq \{ Replace(S) ~|~ S \in \Sigma \} \]
% $Replace(\Sigma) \triangleq \{ Replace(S) ~|~ S \in \Sigma \}$
Similarly, we denote the $GetMinimal$ function:
\[ GetMinimal(\Sigma) \triangleq \{ GetMinimal(S) ~|~ S \in \Sigma \} \]
% $GetMinimal(\Sigma) \triangleq \{ GetMinimal(S) ~|~ S \in \Sigma \}$

\begin{deff}
Given a smart contract, the set of {\em all possible view materialization plans} obtained by applying the Replace function for $k$ times is defined as:
       \[ \Sigma_k \triangleq \bigcup_{i=0}^{k} Replace^i(BasePlan) \] 
\end{deff}

Here, $Replace^i(BasePlan)$ denotes the set of view materialization plans obtained by applying the $Replace$ function (Algorithm~\ref{alg:enumerate}, {\sf replace} function) to the base materialization plan $i$ times. Since circular dependencies among relations are prohibited, the corresponding replacement graph $G$ has no circle. Given $k$ as the length of the longest path in $G$, the $Replace$ function can be applied at most $k$ times.

\begin{deff}\label{def:min-replace}
We define $MinReplace$ function as the combination of the $Replace$ and $GetMinimal$ functions performed at every iteration of Algorithm~\ref{alg:enumerate}. Assuming $\Sigma$ represents a set of view materialization plans, we define $MinReplace$ as follows:
        \[ MinReplace(\Sigma) \triangleq GetMinimal(Replace(\Sigma))  \]
\end{deff}
        
With this definition, the output of Algorithm~\ref{alg:enumerate} is effectively:
\[  \bigcup_{i=0}^{K}MinReplace^i(B') \]
where $B'$ is the minimal form of the base view materialization plan.

\begin{lmm}\label{lmm:getmini}
For any valid materialization plan $S$, 
applying $MinReplace$ method to $S$ is equivalent to first applying $GetMinimal$ to $S$ and then applying $MinReplace$:
\[
    MinReplace(\{S\}) = MinReplace(GetMinimal(\{S\}))
\]
\end{lmm}

This lemma can be further developed to fit a set of valid materialization plans,
instead of just one plan.

\begin{proof}
We first prove $MinReplace(\{S'\}) \subseteq MinReplace(\{S\})$. $\forall$ relation $r \in S$, (1) if $r \in S'$, then $r$ can be executed by any replacement choices of $S'$ including the reduced one $v$; (2) if $r \notin S'$, then $r$ must be calculated using some relations $p \in P$ where $P \subseteq S'$. In this case, $\forall p \in P$, we have $p \in S'$, and therefore $p$ can be executed by any replacement choices of $S'$ including the reduced one $v$, and thus $r$ can also be executed by $v$. Therefore, $v$ is also a reduced alternative choice of $S$. 
We next prove $MinReplace(\{S\}) \subseteq MinReplace(\{S'\})$.
$\forall$ relation $r \in S'$, since $S' \subseteq S$, we have $r \in S$, and therefore $r$ can be executed by any replacement choices of $S$ including the reduced one $u$. So $u$ is also a reduced alternative choice of $S'$.
\end{proof}

\begin{thm}\label{theorem:completeness}{Completeness.}
Given a set of relational queries, let $B$ be the base materialization plan, and $B'$ be the minimal form of $B$ ($GetMinimal(B)$),
Algorithm~\ref{alg:enumerate} outputs all minimal forms of view materialization plans obtainable from $\Sigma$:
        \[ \bigcup_{i=0}^{K}MinReplace^i(B') =  GetMinimal(\Sigma_K) \]
\end{thm}

This theorem asserts that the application of the $GetMinimal$ function during each iteration of the algorithm does not compromise the final correctness of the algorithm. Instead, it enhances efficiency by early pruning of non-minimal materialization plans.
    

\begin{proof}
Given Lemma~\ref {lmm:getmini}, we can prove completeness by induction on the $MinReplace$ step. 
Essentially, at each iteration $k$, the application of $MinReplace$ function maintains completeness regarding the view materialization plan, 
considering $k$ applications of the $replace$ function on the base materialization plan $\Sigma_k$.

\noindent {\bf Base case.} The base case, where $i = 0$, and $B' = GetMinimal(B)$, is true following our assumption.

\noindent {\bf Inductive hypothesis.} The induction case where $i=n$, the inductive hypothesis is:
\[ \bigcup_{i=0}^{n}MinReplace^i(B') =  \{ GetMinimal(S)  | S \in \Sigma_n \} \]


\noindent {\bf Inductive step.} We then analyze $i=n+1$, the output can be rewritten, and further substituted with the induction hypothesis:

\begin{align*}
    \bigcup_{i=0}^{n+1} & MinReplace^i(B') \\
    & = \bigcup_{i=0}^{n}MinReplace^i(B') \cup minReplace^{n+1}(B') \\
    & = \{ GetMinimal(S)  | S \in \Sigma_n \}
                                           \cup  minReplace^{n+1}(B') \\
\end{align*}

Next, we expand $minReplace^{n+1}(B')$: 
\begin{align*}
    minReplace^{n+1}(B') & = minReplace^n(minReplace(B')) \\
\end{align*}

By Definition~\ref{def:min-replace} and lemma~\ref{lmm:getmini}, we can further rewrite the right-hand-side: 
\begin{align*}
    minReplace^{n+1}(B') & = minReplace^n(minReplace(B')) \\
                    & = minReplace^n(minReplace(GetMinimal(B))) \\
                    & = minReplace^n(minReplace(B)) \\
                    & = minReplace^n(GetMinimal(Replace(B))) \\
                    & = minReplace^n(Replace(B)) \\
                    & ... \\
                    & = minReplace(Replace^{n}(B)) \\
                    & = GetMinimal(Replace^{n+1}(B)) \\
                    & = \{ GetMinimal(S) ~|~ S \in Replace^{n+1}(B) \} \\
\end{align*}

Thus, by substituting the equivalent expansion of $minReplace^{n+1}(B')$ to the inductive equation,
we have:
\[
    \bigcup_{i=0}^{n+1}MinReplace^i(B') = \{ GetMinimal(S) ~|~ S \in \Sigma_{n+1} \}
\]
which establishes the validity for $i=n+1$.

Here, we have successfully completed the induction step and prove the completeness property
as defined in Theorem~\ref{theorem:completeness}.
\end{proof}
